% Fichier généré automatiquement par tools/generate_corrections.py.
% Source : SLCI/SLCI-03-SchemaBlocs/77_ProtheseTibia/77_ProtheseTibia.tex
% Statut : complet (3/3 question(s) renseignée(s), 0 reprise(s)).
\begin{corrigebox}[Corrigé extrait de la banque]
\CorrigeQuestion{1}
On a d'une part, $C_M(p)=H_1(p)\left(U_M(p)-\Omega_M(p)\right)$. 

D'autre part, en utilisant les deux équations du moteur électrique, on a $U_M (p)=RI (p)+E(p)$ et $E (p )=k_c \Omega_M (p)$ soit  $U_M (p)=RI (p)+k_c \Omega_M (p)$. De plus $C_M (p)=k_c I (p )$; donc $U_M (p)=R\dfrac{C_M(p)}{k_c}+k_c \Omega_M (p)$. 
Par suite, $C_M(p)=\dfrac{k_c}{R}\left(U_M (p)-k_c \Omega_M (p)\right)$.

En identifiant, on a donc $H_1(p)=\dfrac{k_c}{R}$ et  $H_6(p)={k_c}$.

D'après le schéma-blocs, 

$\Delta \alpha(p) = \left(C(p)-C_M(p)H_2(p)\right)H_3(p)H_4(p)$ soit 

En utilisant l'équation différentielle caractéristique du comportement de la prothèse, on a : 
$J_M p^2 \Delta \alpha(p) + \mu_m p \Delta \alpha(p)  = C_M(p)R_T -C(p)R_T^2$
$\Leftrightarrow \Delta \alpha(p) \left(J_M p^2  + \mu_m p \right)  = C_M(p)R_T -C(p)R_T^2$

$\Leftrightarrow \Delta \alpha(p)  = \dfrac{R_T^2}{J_M p^2  + \mu_m p}\left(\dfrac{C_M(p)}{R_T} -C(p)\right)$.

Or, $\Delta \alpha(p) = \dfrac{1}{p}\Delta \alpha'(p)$; donc $H_4(p)=\dfrac{1}{p}$.

Au final, $H_3(p) = \dfrac{R_T^2}{J_M p  + \mu_m }$ et $H_2(p)=R_T$.
\CorrigeQuestion{2}
On déplace le dernier point de prélèvement avant $H_4$. On ajoute donc $H_4(p)H_7(p)$ dans la retour. 

On a alors $F(p)=\dfrac{\Delta \alpha'(p)}{-} = \dfrac{H_3(p)}{1+H_3(p)H_4(p)H_7(p)}$.
$\text{FTBF}(p)  = \dfrac{H_1(p)H_2(p) F(p)}{1+H_1(p)H_2(p)H_5(p)H_6(p) F(p)}H_4(p)H_7(p) $.

Soit 
$\text{FTBF}(p)  = \dfrac{H_1(p)H_2(p) \dfrac{H_3(p)}{1+H_3(p)H_4(p)H_7(p)}}{1+H_1(p)H_2(p)H_5(p)H_6(p) \dfrac{H_3(p)}{1+H_3(p)H_4(p)H_7(p)}}H_4(p)H_7(p) $

$= \dfrac{H_1(p)H_2(p) H_3(p)}{1+H_3(p)H_4(p)H_7(p)+H_1(p)H_2(p)H_5(p)H_6(p) H_3(p)}H_4(p)H_7(p) $

$= \dfrac{\dfrac{k_c}{R}R_T \dfrac{R_T^2}{J_M p  + \mu_m }}{1+\dfrac{R_T^2}{J_M p  + \mu_m }\dfrac{k_{RS}d_0^2}{p}+\dfrac{k_c}{R}R_T \dfrac{1}{R_T}k_c \dfrac{R_T^2}{J_M p  + \mu_m }}\dfrac{k_{RS}d_0^2}{p} $

$= \dfrac{\dfrac{k_c}{1} R_T^3}{J_MR p^2  + \mu_m R p+R_TR^2k_{RS}d_0^2+pk_ck_c R_T^2}k_{RS}d_0^2 $

$= \dfrac{k_c R_T^3}{J_MR p^2  + p\left(\mu_m R  +k_ck_c R_T^2\right)+R_TR^2k_{RS}d_0^2}k_{RS}d_0^2 $.
\CorrigeQuestion{3}
\begin{itemize}
\item Le régime permanent semble atteint autour de \SI{0,03}{s}; donc les critère de rapidité est respécté.
\item En régime permanent, le couple atteint est de \SI{46}{Nm} pour une consigne de \SI{50}{Nm}. Un écart de 10 \% correspondrait à un couple atteint de \SI{45}{Nm}. Le critère de précision est respecté.
\end{itemize}

\footnotesize
\begin{enumerate}
\item  $H_1(p)=\dfrac{k_c}{R}$, $H_2(p)=R_T$,  $H_3(p) = \dfrac{R_T^2}{J_M p  + \mu_m }$ et   $H_6(p)={k_c}$.
\item $\text{FTBF}(p)= \dfrac{k_c R_T^3}{J_MR p^2  + p\left(\mu_m R  +k_ck_c R_T^2\right)+R_TR^2k_{RS}d_0^2}k_{RS}d_0^2 $.
\item .
\end{enumerate} 
\normalsize
\end{corrigebox}
